\chapter{数学符号与公式索引}
\label{app:symbols-formulas}

本附录只索引全书反复使用的符号；局部变量以首次出现处定义为准。

\begin{longtable}{p{0.16\textwidth}p{0.48\textwidth}p{0.22\textwidth}}
\caption{核心符号索引}\label{tab:core-symbols}\\
\toprule
符号 & 含义 & 主讲章节 \\
\midrule\endfirsthead
\toprule 符号 & 含义 & 主讲章节 \\
\midrule\endhead
$\mathbf{x},\mathbf{u}$ & 状态与控制/动作向量 & 第 4、7、10、22 章 \\
$\mathbf{T}\in SE(3)$ & 三维刚体位姿齐次变换 & 第 5 章 \\
$\boldsymbol{\xi}\in\mathfrak{se}(3)$ & 位姿李代数坐标/扭量 & 第 5 章 \\
$\mathbf{J}(\mathbf{q})$ & 关节速度到末端速度的雅可比 & 第 5、6 章 \\
$\mathbf{M},\mathbf{C},\mathbf{g}$ & 质量矩阵、科氏/离心项与重力项 & 第 6 章 \\
$\boldsymbol{\tau},\boldsymbol{\lambda}$ & 关节力矩与接触约束力 & 第 6 章 \\
$\pi(\mathbf{a}\mid\mathbf{o})$ & 观测条件下的策略分布 & 第 9--12、22 章 \\
$\gamma,\mathcal{R}$ & 折扣因子与回报/奖励 & 第 10 章 \\
$Q,K,V$ & Transformer 查询、键和值 & 第 14 章 \\
$H_s,A,r,Y$ & 排程时间、可用率、产率与良率 & 第 31、35 章 \\
$C_{\mathrm{TCO}}$ & 合同期总拥有成本 & 第 35 章 \\
\bottomrule
\end{longtable}

公式阅读顺序是：假设与坐标系、变量与单位、推导关系、数值算例、失效条件。只引用公式编号而脱离这些上下文会造成误用。

